%% ============================================================
%%  SECTION 6 — KẾT LUẬN (Conclusion)  (số liệu: v2)
%% ============================================================
\section{Kết luận}
\label{sec:conclusion}

\subsection{Tóm tắt quá trình và kết quả}

Đồ án đã hoàn thành trọn vẹn quy trình Khoa học Dữ liệu (KDD) 6 bước:
thu thập 133.586 bản ghi từ 10 nguồn, tích hợp về lược đồ thống nhất,
làm sạch và phân tích khám phá, xây dựng và đánh giá 5 mô hình hồi quy,
kết thúc bằng ứng dụng demo Streamlit.

Mô hình triển khai \textbf{XGBoost} đạt:
\begin{itemize}[leftmargin=*, itemsep=2pt]
  \item \textbf{Test R\textsuperscript{2} = 0.565}; \textbf{Test RMSE = 1.895 triệu/tháng};
        \textbf{Test MAE = 1.26 triệu/tháng}.
  \item So sánh công bằng (cùng ràng buộc đơn điệu cho cả hai mô hình cây): XGBoost
        \textbf{chính xác hơn} Random Forest (Test R\textsuperscript{2} $0.565$ vs.\ $0.459$) —
        \textit{không có} đánh đổi độ chính xác lấy khả năng triển khai.
  \item Kích thước $\approx$0.86 MB, triển khai được trong giới hạn GitHub và app Streamlit.
\end{itemize}

\subsection{Bài học kinh nghiệm và ý nghĩa}

\begin{enumerate}[leftmargin=*, itemsep=5pt]
  \item \textbf{Diện tích và khu vực là yếu tố định giá mạnh nhất.}
        Feature Importance của cả RF lẫn XGBoost đều xếp \texttt{area}
        và \texttt{city}/\texttt{district} dẫn đầu sau khi loại \texttt{source}.

  \item \textbf{Nghịch lý Simpson là bài học về phân tích dữ liệu nghiêm túc.}
        Tương quan âm toàn cục giữa tiện ích và giá không phản ánh quan hệ nhân quả —
        đây là confounding variable do cấu trúc nguồn.
        Phát hiện này dẫn tới quyết định loại \texttt{source}
        (đánh đổi một phần độ chính xác nhưng đổi lại mô hình \textbf{trung thực và tổng quát hơn}).

  \item \textbf{So sánh mô hình phải công bằng — áp cùng điều kiện cho mọi ứng viên.}
        Khi chỉ XGBoost bị ràng buộc đơn điệu còn Random Forest chạy tự do, RF có
        R\textsuperscript{2} cao hơn — nhưng đó là so sánh khập khiễng: RF ``thắng'' nhờ tự do khai
        thác đúng những tương quan tiện ích phi lý mà ràng buộc cố ý loại bỏ. Sau khi áp
        \textbf{cùng} ràng buộc đơn điệu cho cả hai, XGBoost vượt Random Forest
        ($0.565$ vs.\ $0.459$): nó vừa chính xác nhất vừa gọn nhẹ ($\approx$0.86 MB so với
        $\approx$56 MB), suy luận nhanh và có early stopping. Lựa chọn mô hình triển khai vẫn phải
        cân nhắc cả \textbf{ràng buộc kỹ thuật} (kích thước, tốc độ, bảo trì), không chỉ độ chính xác.

  \item \textbf{Các quyết định tiền xử lý phải được kiểm chứng bằng thực nghiệm, không chỉ lập luận.}
        Ba ablation study xác nhận:
        \begin{itemize}[itemsep=2pt]
          \item FE (tách khu vực từ địa chỉ) giảm RMSE 3.5--14.6\%;
                PCA không giúp ích với bộ dữ liệu này.
          \item Log-transform mục tiêu \textit{không} tăng hiệu năng theo tiêu chí RMSE/R\textsuperscript{2}
                và gây underweight phân khúc giá cao.
          \item Lọc ngoại lệ giảm RMSE 79.0\% trên vùng giá mục tiêu.
        \end{itemize}

  \item \textbf{Quan hệ phi tuyến mạnh.}
        Mô hình cây (R\textsuperscript{2} $\approx$ 0.57--0.66) vượt xa tuyến tính
        (R\textsuperscript{2} $\approx$ 0.19), xác nhận tương tác
        khu vực $\times$ diện tích không thể biểu diễn bằng hàm tuyến tính đơn giản.

  \item \textbf{Giới hạn bởi bias (không phải variance).}
        Learning curve cho thấy thêm dữ liệu cùng loại giúp ít —
        cần đặc trưng \textit{mới} (GPS, khoảng cách, NLP mô tả) để tăng R\textsuperscript{2}.
\end{enumerate}

\subsection{Hướng phát triển (Future Work)}

\begin{itemize}[leftmargin=*, itemsep=3pt]
  \item \textbf{Bổ sung và chuẩn hóa dữ liệu tiện ích} đồng đều giữa các nguồn để
        tăng chất lượng đặc trưng nhị phân, từ đó cải thiện độ chính xác mô hình.
  \item \textbf{Đặc trưng vị trí không gian:} tọa độ GPS, khoảng cách đến trung tâm,
        tiện ích công cộng (trường học, bệnh viện, metro) — các nghiên cứu quốc tế
        cho thấy đây là yếu tố tăng R\textsuperscript{2} mạnh nhất.
  \item \textbf{Khai thác mô tả dạng văn bản:} NLP (TF-IDF, BERT) trên trường
        \texttt{description} để trích xuất tín hiệu giá (tầng, hướng, view, trang trí…).
  \item \textbf{Cân bằng dữ liệu địa lý:} thu thập thêm mẫu ngoài Hà Nội
        để mô hình đại diện tốt hơn cho các tỉnh thành khác.
  \item \textbf{Cập nhật định kỳ:} triển khai pipeline tự động cào và retrain
        hàng tháng, theo dõi data drift khi thị trường thay đổi.
  \item \textbf{Mô hình phức tạp hơn:} LightGBM, CatBoost, hoặc deep learning
        (TabNet) nếu bổ sung được đặc trưng GPS và văn bản.
\end{itemize}
